\section{讨论}
\label{zh:sec:discussion}

\subsection{证据能够确立的内容}

本文的主要结果是约束接口效应，而不是声称某一个优化器支配所有替代方法。MIKU 保留 PVD 求解器，但改变障碍物信息何时以及以何种方式到达求解器。最大收益出现在窄间隙多障碍物和延迟交叉构型中，恰好是预测时域并集或不可逆局部侧向选择会移除可行替代方案的情形。在 B0 已经容易处理的构型中，例如交叉行人和车辆切入类别，总体成功率没有变化。相比所有案例都统一提升，这种结果模式更直接地支持本文提出的机制。

消融实验提供了第二层支持。Top-$K$ 空间同伦和威胁感知裕度解释了总体成功率差异的大部分，而时间同伦图主要在障碍物前后两种通行均可到达时发挥作用。鲁棒占用管在预测噪声下减少碰撞，但相应的进度和加加速度代价表明，面向安全的扩展会消耗可用空间。这些权衡与将可行性、安全性和舒适性视为竞争性运营结果的走廊规划和预测感知规划研究一致 \citep{yoon2024stcorridor,yangbo2023ped,pathspeedstructured2024}。

\subsection{运营含义}

对于基于 PVD 的自动驾驶车辆规划栈，一个可操作的设计规则是在确定路径边界之前保留交互时序，然后将选定的时间语义传递给速度阶段。这样可以在不引入新的下游非线性或混合整数求解器的情况下，减少不必要的兜底停车。该规则尤其适用于窄施工区、临时车道绕行以及多个障碍物共享同一纵向分量的混合交通场景。部署时仍应配备延迟监控器和兜底策略，因为运行时间上尾和不确定性扩展案例仍然具有实际影响。

\subsection{与替代方法的关系}

有限网格参考在小规模比较中取得略高的成功率，但运行时间中位数约高出两个数量级。这一比较说明扩展联合搜索与编译式 PVD 接口之间存在质量—代价权衡，并不能证明 MIKU 在全局上更优。同样，迭代式 B2 基线表明，当缺失信息是已经丢弃的横向类别时，简单重复 PVD 阶段并不足够。因此，MIKU 的贡献在于约束的表示和传递，而不是声称所有部署都不需要迭代或联合搜索。

\subsection{对评估的启示}

3,500 个场景的数值协议、700 个场景的滚动协议和 Apollo 接口案例彼此分开，对于正确解释结果十分重要。第一类估计匹配场景层面的效应，第二类测试重规划下的累积行为，第三类检查工程插入点。将它们合并为一个总分会夸大外部有效性。后续评估应使用公共或道路采集轨迹、与近期外部规划器的匹配比较，并在目标硬件预算下进行原生运行时间测量。
